\documentclass[output=paper,colorlinks,citecolor=brown]{langscibook}
\ChapterDOI{10.5281/zenodo.20541844}


\author{Lena Schnee\orcid{}\affiliation{Universität Hamburg} and
        Stefan Hartmann\orcid{}\affiliation{Heinrich-Heine-Universität Düsseldorf}}
\title{Introduction}
\abstract{This introduction gives a very brief overview of the state of the art in research on the linguistic expression of future reference, as well as on the diachronic development of future constructions across languages. It also offers a brief summary of the individual chapters to the present edited volume, discussing their contributions to the overarching research questions addressed in this book, and outlines potential avenues for further contrastive research on the diachrony of future constructions.}



\IfFileExists{../localcommands.tex}{
   \addbibresource{../localbibliography.bib}
   \usepackage{tabularx,multicol}
\usepackage{url}
\urlstyle{same}

\usepackage[english]{babel}
 
\usepackage{langsci-optional}
\usepackage{langsci-lgr}
\usepackage{langsci-gb4e}

\usepackage{langsci-branding}

%\usepackage{tabularx}

% elena
\usepackage{graphicx} %package to manage images
\usepackage{paralist}

% sarah & lena
\usepackage[justification=centering]{caption}
\usepackage{rotating}

% robin
\usepackage{lscape}
\usepackage{tikz}
\usetikzlibrary{tikzmark}

% dominika
\usepackage{multirow}
\usepackage{hanging}

% olaf & dylan
\usepackage{subcaption}

% svetlana
\usepackage{verbatim}
\usepackage{enumitem}

% wiemer
%\usepackage[markup=underlined]{changes}
%\usepackage{hyphenat}

% wir
\usepackage{longtable}

% for crossref in the introduction jiabin
\usepackage{xr}
   % \newcommand*{\orcid}{}
%

\makeatletter
\let\thetitle\@title
\let\theauthor\@author
\makeatother

\newcommand{\togglepaper}[1][0]{
   \bibliography{../localbibliography}
   \papernote{\scriptsize\normalfont
     \theauthor.
     \titleTemp.
     To appear in:
     E. Di Tor \& Herr Rausgeberin (ed.).
     Booktitle in localcommands.tex.
     Berlin: Language Science Press. [preliminary page numbering]
   }
   \pagenumbering{roman}
   \setcounter{chapter}{#1}
   \addtocounter{chapter}{-1}
}
%
\newbool{bookcompile}
\booltrue{bookcompile}
\newcommand{\bookorchapter}[2]{\ifbool{bookcompile}{#1}{#2}}
%
% %\newcommand{\orcid}[1]{}
\graphicspath{ {./figures/} }
%
\let\eachwordone=\itshape


% Seongha
% \newcommand{\koreex}[1]{{\fontspec{Baekmuk Batang}{#1}}}

% Kutida
\newfontfamily\thaifont{Noto Sans Thai} 
\newcommand{\textthai}[1]{{\thaifont #1}} % für Thai


\newfontfamily\chinesefont{Noto Sans CJK SC}
\newcommand{\textchinese}[1]{{\chinesefont #1}} % Für Chinesisch

% wiemer
\def\hyph{-\penalty0\hskip0pt\relax}
% \definechangesauthor[color=NavyBlue]{SH}



\newcolumntype{K}{>{\raggedright\arraybackslash}p{0.6\textwidth}}  % Right column with 0.7 width


% \newfontfamily\krn[Scale=MatchLowercase,BoldFont=SourceHanSerif-Bold.otf]{SourceHanSerif-Regular.otf}
% \AdditionalFontImprint{Source Han Serif KO}
% \XeTeXlinebreaklocale 'ko'
% \renewcommand{\krn}{}



% Rhee
\newfontfamily\koreanfont{Noto Sans CJK KR} 
\newcommand{\textkorean}[1]{{\koreanfont #1}} % für Koreanisch

\renewcommand{\REF}[2][]{(\ref{#2}#1)}
   %% hyphenation points for line breaks
%% Normally, automatic hyphenation in LaTeX is very good
%% If a word is mis-hyphenated, add it to this file
%%
%% add information to TeX file before \begin{document} with:
%% %% hyphenation points for line breaks
%% Normally, automatic hyphenation in LaTeX is very good
%% If a word is mis-hyphenated, add it to this file
%%
%% add information to TeX file before \begin{document} with:
%% %% hyphenation points for line breaks
%% Normally, automatic hyphenation in LaTeX is very good
%% If a word is mis-hyphenated, add it to this file
%%
%% add information to TeX file before \begin{document} with:
%% \include{localhyphenation}

\hyphenation{
    con-tri-bute
    par-a-digm
    pe-ri-phra-ses
    achieve-ment
    accomp-lish-ment
    We-li-nchen-kang-ci-kok
    cross-ling-u-is-tic
    rec-ord-s
    trun-ca-ted
    de-ri-va-tion-al
    stem-de-ri-va-tion-al
    pre-sent
    in-flu-en-ced
    sys-tem
    ques-tions
    Bor-kovs-kij
    Da-ny-len-ko
    e-qui-va-lents
    re-cog-nized
    In-ter-fe-renz-schicht
    Aus-ei-nan-der-se-tz-ung
    Früh-neu-hoch-deutsch-kor-pus
    Ak-tions-art
    Smir-no-va
    Whi-le
    Ar-chive
    bi-ginnan
    Rei-chen-bach
    ger-ma-nis-ti-sche
    Re-fe-renz-kor-pus
    le-xi-ka-lisch-en
    Alt-hoch-deutsche
    Li-te-ra-tur
    dia-chron
    Dia-lek-to-lo-gie
}

\hyphenation{
    con-tri-bute
    par-a-digm
    pe-ri-phra-ses
    achieve-ment
    accomp-lish-ment
    We-li-nchen-kang-ci-kok
    cross-ling-u-is-tic
    rec-ord-s
    trun-ca-ted
    de-ri-va-tion-al
    stem-de-ri-va-tion-al
    pre-sent
    in-flu-en-ced
    sys-tem
    ques-tions
    Bor-kovs-kij
    Da-ny-len-ko
    e-qui-va-lents
    re-cog-nized
    In-ter-fe-renz-schicht
    Aus-ei-nan-der-se-tz-ung
    Früh-neu-hoch-deutsch-kor-pus
    Ak-tions-art
    Smir-no-va
    Whi-le
    Ar-chive
    bi-ginnan
    Rei-chen-bach
    ger-ma-nis-ti-sche
    Re-fe-renz-kor-pus
    le-xi-ka-lisch-en
    Alt-hoch-deutsche
    Li-te-ra-tur
    dia-chron
    Dia-lek-to-lo-gie
}

\hyphenation{
    con-tri-bute
    par-a-digm
    pe-ri-phra-ses
    achieve-ment
    accomp-lish-ment
    We-li-nchen-kang-ci-kok
    cross-ling-u-is-tic
    rec-ord-s
    trun-ca-ted
    de-ri-va-tion-al
    stem-de-ri-va-tion-al
    pre-sent
    in-flu-en-ced
    sys-tem
    ques-tions
    Bor-kovs-kij
    Da-ny-len-ko
    e-qui-va-lents
    re-cog-nized
    In-ter-fe-renz-schicht
    Aus-ei-nan-der-se-tz-ung
    Früh-neu-hoch-deutsch-kor-pus
    Ak-tions-art
    Smir-no-va
    Whi-le
    Ar-chive
    bi-ginnan
    Rei-chen-bach
    ger-ma-nis-ti-sche
    Re-fe-renz-kor-pus
    le-xi-ka-lisch-en
    Alt-hoch-deutsche
    Li-te-ra-tur
    dia-chron
    Dia-lek-to-lo-gie
}
   \boolfalse{bookcompile}
   \togglepaper[1]%%chapternumber
}{}

\begin{document}
\maketitle

\noindent Thinking and talking about future events is part of our everyday lives. As such, it is not surprising that languages have developed a variety of markers for futurity. According to \citet{wals-67}, ``[i]t is relatively rare for a language to totally lack any grammatical means for marking the future. Most languages have at least one or more weakly grammaticalized devices for doing so." In \citegen{BybeePerkinsPagliuca1994} broad typological sample, 70 out of 90 languages had at least one grammatical marker for future time reference, and among these, 49 had more than one form to refer to the future \citep[also see][2]{wiemer_between_2024}. The fact that many languages have more than one way of expressing future time reference can be seen as a reflection of the diachronic tendency of futures to be unstable over time and to be formed anew in polygenesis \citep[][2]{wiemer_between_2024}. This instability makes them an even more interesting and promising object of study from a diachronic and contrastive point of view. Crucially, there are vast typological differences between languages, both in the degree to which future markers are grammaticalized and in the diachronic paths of development that these markers have taken. The present volume focuses on the latter aspect: Its goal is to offer in-depth studies on the diachronic development of future markers across different languages. 

In doing so, it connects to a rich research tradition that has addressed the evolution of grammatical markers for future time reference, which in turn is part of a broader literature on tense and aspect. Fleischman (\citeyear{Fleischman1982}: 8,  \citeyear{fleischman_tense_1990}: 26) describes the temporal structure of future tenses as either S>R/E or S/R>E, taking up \citegen{Reichenbach1947} tense model that distinguishes between the point in time when an utterance is made (point of speech, S), the point in time when an event occurs (point of the event, E), and the point of reference (R), which can coincide with S or E but can also be independent, e.g. in the case of the future perfect (\textit{I will have seen John}: S > E > R; \citeauthor{Reichenbach1947} \citeyear{Reichenbach1947}: 3).  In his broad typological overview of future tenses, \citet[]{Ultan1978} argues that most languages show either a prospective or a retrospective tense system: in the former, the present tense can generally be used to refer to the future; in the latter, the present tense can generally be used to refer to the past. Regardless of whether a language has a prospective or a retrospective tense system, he argues that ``future tenses are generally more marked than either past or present tenses" \citep[][116]{Ultan1978}, which is reflected in the markedness of the grammatical features of future tenses. For instance, if a language has a periphrastic past tense (which he considers more marked than bound suffixes used for tense marking), it will also have a periphrastic future, but not vice versa \citep[][90--91]{Ultan1978}.

\citet[309--310]{Dahl2000} points out that the future has a very different epistemological status compared to the present or the past: ``We cannot perceive or remember future states of affairs, and it has been disputed whether statements about the future can be said to have a determinate truth value." However, we do talk a lot about intentions on the one hand and predictions on the other, which leads \citet[310]{Dahl2000} to the distinction between prediction-based and intention-based future time reference. He observes that in European languages, markers that are originally restricted to intention-based future time reference tend to develop into general future makers; at the same time, a grammatical opposition between intention- and prediction-based future time reference ``is less common than one would perhaps
think in view of the apparent cognitive salience of that distinction" \citep[][310]{Dahl2000}.

\citet[244]{BybeePerkinsPagliuca1994} propose a distinction between primary futures, which evolve from a fairly restricted set of lexical sources (constructions involving movement verbs; markers of obligation, desire, and ability; temporal adverbs) on the one hand and what they call aspectual futures on the other, which ``use a form whose principal function is the marking of present tense or perfective or imperfective aspect". The latter category would include, for example, cases of futurate present such as English \textit{The Yankees play the Red Sox tomorrow.} \citep[][158]{Hilpert2008}, which would not typically be considered an instance of a future construction, let alone a future \textit{tense}. This is why some terminological remarks are in order. When we talk about future \textit{constructions} in the subtitle of the present book, we mean linguistic patterns that can be used for future time reference. We use the term in a fairly theory-neutral sense here, although it is clearly inspired by the notion of constructions in the sense of form-meaning pairs in constructionist approaches to language \citep[][]{ungerer_constructionist_2023}. Note, however, that the term is used in different senses in the literature, including the contributions to the present volume. For example, \citet[1]{wiemer_between_2024} use it interchangeably with \textit{marker} or \textit{gram}. Especially from a diachronic point of view, the distinction between a pattern that is a fully-fledged grammatical marker for future time reference and one that can be \textit{used} for referring to future event even though this is not its primary purpose is of course blurred, which is why the scope of the present volume is not limited to cases that are clearly identifiable as future tenses -- in fact, there have been controversial debates whether the German \textit{werden} + infinitive construction, which several of the papers deal with, is actually a future construction or not \citep[e.g.][]{Saltveit1962, vater1975}. What counts for our purposes is, first and foremost, that a pattern is used (more or less regularly) for future time reference. 

The present volume does not aim to cover the full typological breadth of languages, or of grammaticalization scenarios of future constructions. Instead, it presents some case studies from a diverse set of languages that can be considered representative for some of the recurrent questions that are discussed with regard to the diachronic development of linguistic markers of future reference. 

A contrastive perspective on the diachrony of future constructions can prove insightful for a variety of reasons. Firstly, the developmental patterns of future constructions constitute prime examples of grammaticalization \citep{HopperTraugott2003}. But future constructions differ in their grammaticalization paths:  \linebreak While some derive from modal verbs (e.g. English \textit{will}/\textit{shall}), others recruit \linebreak change-of-state verbs (e.g. German \textit{werden} `become') or forms of \textit{come} and \textit{go} (e.g. English \textit{BE going to}; Duala -\textit{ya} ‘go’, Bambara \textit{nà} ‘come’, among many others, see \cite{Kuteva2019}) as future auxiliaries. Secondly, the competition patterns between future constructions merit further investigation from a cross-linguistic perspective. The examples mentioned above also show that in many cases, future constructions going back to verbs of movement compete with future constructions that have developed from modal verbs \citep[see e.g.][]{Kuteva2019}. In other languages, change-of-state verbs such as German \textit{werden} ‘become’ have prevailed as future markers \citep[][]{Dahl1985, Dahl2000}. Crucially, languages differ not only in the ways in which they mark future time reference, but also in the extent to which speakers can choose between future tense forms and the so-called futurate present when referring to future events \citep[see e.g.][]{wals-67}.

A third aspect is equally important: The diachronic development of future constructions, as well as their synchronic diversity, can potentially give important clues to our cognitive conceptualizations of time, as well as to cultural factors that may play a role in how time is conceptualized and expressed linguistically \citep{Fleischman1982}. In fact, \citet[]{coseriu1957} puts forward the far-reaching hypothesis that the rise of Christianity was a major factor in the emergence of modal futures in Romance languages. While this idea has to be taken with some scepticism (see \citeauthor{Fleischman1982} \citeyear{Fleischman1982}: 49--50 for critical remarks), it does seem plausible to assume that socio-historical and cultural phenomena can have an influence on the use of grammatical forms. In addition, the influence of writing and literacy should probably not be underestimated, as differences between oral and literal contexts have often been found to be crucial for tense uses, especially when it comes to the use of future constructions \citep[see e.g.][]{fleischman_tense_1990, nau_future_2022}.

The goal of this volume is thus to take a contrastive perspective on the evolution of future constructions by bringing together scholars working on different languages and language families. In particular, we want to address the following questions:

 \begin{itemize}
     \item Which cross-linguistic tendencies can be observed in the grammaticalization of future constructions?
     \item To what extent can we compare the grammaticalization paths of future constructions with different origins?
     \item How do different future constructions compete in different languages, and can these patterns of competition be compared cross-linguistically?
     \item Which cognitive implications do findings on the diachronic development of future constructions entail?
 \end{itemize}

As for the first question, a number of contributions re-address common hypotheses about the main grammaticalization paths of future constructions, while others, like Rhee and Khammee, delve deeper into the intricacies of grammaticalization paths of future constructions in two languages, Thai and Korean. In addition, grammaticalization processes can be intertwined with other processes of language change that are not traditionally considered part of grammaticalization. For instance, hypoanalysis is identified as a key process with respect to both the development of the Middle Armenian Future (Meyer) as well as the futurate present of perfective Verbs in North Slavic languages (Wiemer). These contributions also address the third question mentioned above, i.e. parallel patterns of competition between different future constructions in different languages. 

Importantly, the contributions to our volume also address languages and language stages for which no systematic analysis of future tenses exists so far, e.g. Middle Low German (Ihden \& Schnee). They also address understudied topics such as the encoding of ‘future-in-the-past’ (Petrova \& Chen) or ``failed" grammaticalization paths in the development of future constructions (Skrzypek, Smirnova), which arguably have implications for the fourth question raised above, i.e. for a deeper understanding of the cognitive mechanisms underlying the grammaticalization of future reference markers.

Another aspect that the contributions to this volume add to the existing literature is an in-depth discussion of data types that are suitable for investigating the development of future markers. Bible translations are a widely-used resource in this domain as they are often based on the Latin version of the Bible, the Vulgata, and as Latin has a morphological future marker, it can be revealing to check how Latin future forms are translated (Ihden \& Schnee). 

In sum, the edited volume opens up new perspectives on long-standing questions regarding the diachronic development of future constructions across languages, with implications not only for research on tense-aspect-mood systems but also for research on grammaticalization and language change more generally. It brings together eleven papers presenting empirical case studies on future constructions in High and Low German, Norwegian, Swedish, West Polesian, Korean, Thai, as well as in multiple Slavic languages. The articles focus on a variety of research questions, some offering a broad-scope overview of previously under-investigated cases of grammaticalization in lesser studied languages, others zooming in on specific aspects pertaining to the development of well-investigated patterns. 

Three papers in the present volume deal with the German future construction \textit{werden} + Infinitive. One of them, by \textbf{Sarah Ihden} and \textbf{Lena Schnee}, focuses on Low German. They present an onomasiological study that investigates the factors that influence the use of different grammatical structures expressing future reference in Middle Low German. In addition to modal verb constructions, which are considered the most common choices in the old Middle Low German grammars, futurate usage of the present tense is also possible, as well as verbal constructions with \textit{werden} with an infinitive or a present participle. Drawing on selected texts from the reference corpus of Middle Low German / Low Rhenisch, Ihden and Schnee check how future constructions are translated in Middle Low German Bible texts, showing that both the translator and the semantic context have a major influence on the choice between different variants. Furthermore, parallels and differences are identified in comparison to the development in High German. 

\textbf{Lena Schnee} and \textbf{Stefan Hartmann} investigate the relation between \textit{werden} + Infinitive and \textit{werden} + Present Participle based on data from historical corpora of German. \textit{werden} + Infinitive is the most important overt future marker of present-day German, but there are many different theories about its emergence and development, in many of which \textit{werden} + Present Participle plays a crucial role. While it is fairly uncontroversial by now that \textit{werden} + Infinitive did not emerge from some kind of reanalysis of \textit{werden} + Present Participle, as assumed in very early approaches (e.g. \textit{werde lachend} `become laughing' > \textit{werde lachen} `will laugh'), it still seems plausible to assume that the two constructions influenced each other in their development. These patterns of competition are investigated using association and dispersion measures. The results provide tentative evidence that over the time span taken into account in the study, the two constructions first become more similar and then start occupying different functional niches, with \textit{werden} + Present Participle largely falling out of use except in a few dialects.

The diachronic change of the German \textit{werden} + Infinitive construction is also the topic of \textbf{Elena Smirnova}'s paper, which investigates its development  in the early stages of New High German. Drawing on data from the German Text Archive and the \textsc{dwds} Core Corpus of the 20\textsuperscript{th} century, she shows that the construction has undergone a decline in frequency, being replaced by the futurate present in many contexts, and that the construction has established a clear preference for the verb \textit{sein} `to be', which indicates that it has come to occupy a fairly narrow functional niche, and that two tendencies that are typically considered characteristic of grammaticalization cannot be observed any more in the case of this pattern: context expansion \citep[][]{Himmelmann2004} on the one hand, and obligatorification \citep[][]{Lehmann2015} on the other. On this account, then, the grammaticalization of \textit{werden} + Infinitive has run into a dead end.

\textbf{Dominika Skrzypek} examines constructions with a \textit{werden}-cognate, \textit{varda}, in Swedish, which did not embark on the future grammaticalisation path after all. Similarly to the situation in German, \textit{varda} was found in a passive construction with the past participle as well as with an infinitive, in future or inchoative contexts. Her corpus study traces the development of \textit{varda} usages from 1300 to 1750, comparing the overall conditions with those in German. Several reasons for the “demise” are discussed, shedding light on the more general question of promoting and hindering factors of grammaticalisation. One main hindering factor for \textit{varda} is the rise of the Middle Low German loan \textit{bliva} ‘remain’, which developed a new meaning ‘become' in Swedish and replaced \textit{varda} in all constructions except the infinitive one, leading to its near disappearance. 

While in Swedish, Middle Low German \textit{bliva} did not start taking an infinitive, it did so in Norwegian: \textbf{Olaf Mikkelsen} and \textbf{Dylan Glynn} investigate the competition between the Norwegian future constructions \textit{blir å} + Infinitive and \textit{vil} + Infinitive. While the former used to be considered a feature of Northern Norwegian dialects, the authors present evidence that it has spread over the last decade, showing considerable functional overlap with \textit{vil }+ Infinitive as well as with other competing future constructions. Mikkelsen and Glynn present a corpus study based on a Norwegian blog corpus, showing that ``\textit{blir å \textsc{inf}} is encroaching on the semantic territory of \textit{vil \textsc{inf}}" and deeming it ``a future that may still be".

Two contributions deal with the relation between grammaticalization and hypoanalysis. Hypoanalysis refers to the phenomenon that a contextual semantic or functional property is reanalyzed as an inherent property of the syntactic unit \citep[see][126]{Croft2000}.  The paper by \textbf{Björn Wiemer}  discusses the interaction of grammaticalization and hypoanalysis in the development of future tenses in Slavic languages. Wiemer shows that in North Slavic languages, the present tense of perfective stems comes to mark perfective future by default. Wiemer argues that this is an instance of hypoanalysis. He slightly adapts Croft's definition of hypoanalysis, allowing for it to be applied not only to cases where a new function was previously only implied but also to cases where a function was always possible but turns into the most prominent one. In South Slavic, by contrast, `become'- and `want'-based futures emerged, which can be seen as results of grammaticalization.

\textbf{Robin Meyer} traces the development of future constructions from Old Armenian to Modern Armenian, presenting a complex scenario and offering a detailed discussion of the relationship between the multiple change mechanisms that have contributed to the development of these constructions.
Old Armenian did not have a morphological future tense, instead, future  events were typically expressed by means of modal forms. While such are a common starting point for the grammaticalisation of future markers,  the Armenian language did not take that pathway. The Middle Armenian future has been shown to have developed from a former present through hypoanalysis and univerbation. In the Eastern varieties of Modern Armenian, additional analytic future periphrases with specific aspectual meanings have arisen, leading to an increase in aspectual differentiation.

\textbf{Svetlana Petrova \& Yen-Chun Chen} address a special domain in the field of future research, the ‘future in the past’ (FiP). They give an account of the means that were used to express this posteriority of an event in a past tense setting in Old High German and Middle High German, prior to the establishment of the \textit{würde}-construction. In their onomasiological corpus study, they find that overall, the past subjunctive correlates most strongly with this function, however they also observe an increase in modal constructions. Comparing the usage patterns of these with the modal constructions found in the canonical future, they find them to behave similarly in both contexts. The diachronic evidence they provide sheds light on controversially discussed issues such as the nature of FiP as a temporal category. In their paper, parallels are outlined between FiP and the canonical future.

The contributions by \textbf{Kultida Khammee} and \textbf{Seongha Rhee} investigate the grammaticalization of future markers in Thai and Korean, respectively. Khammee's contribution shows that there are various ways to mark future time reference in present-day Thai, with \textit{cà}  being the primary future marker; it developed out of \textit{càk}, which also continues to be used as a future marker. The origins of \textit{càk} are unclear but Khammee hypothesizes that it may go back to a  borrowing from Middle Chinese meaning ‘know’, while also outlining a number of other possibilities. In her paper, she discusses possible pathways of grammaticalization of Thai future markers against the background of \citegen{Lehmann2015} grammaticalization parameters. Rhee's paper shows that Korean also has a large number of options for future marking, but unlike many other languages, it has two formants that inherently involve futurity, viz. suffixes with the meaning `(at a) later time'. More specifically, he shows that contemporary Korean has six future markers with varying grammatical status ranging from adnominalizing suffix (\textit{-l}), to a modal auxiliary (-\textit{li-, -l.ke(s).i-, -keyss-}), and sentence-ender (-\textit{li, -l.key}). Similar to Khammee, he discusses to what extent the current status of the future markers can be explained in the light of \citegen{Lehmann2015} grammaticalization parameters. 

The chapter by \textbf{Kristian Roncero} leads us to a completely different language family. In his study of future constructions in West Polesian (East Slavic), he shows that the de-obligative (xoʧu ɾobɪtɪ ‘I will (need to) do’) and de-volitive futures (majusʲ(a) ɾobɪtɪ ‘I am planning to do’) have had a highly unusual course of development in that they have inverted their original semantic function: while the de-volitive has synchronically an obligational nuance, the de-obligative synchronically expresses an intention. While he cannot offer a final explanation of this phenomenon, one potential explanation he proposes is that the overabundance of morphological forms for expressing futurity may have forced some of the patterns to specialize over time; however, he also points to potential additional factors as well as the remaining high degree of variability between the patterns. 

Although the contributions to the present volume cannot represent the full typological breadth and diversity of future marking, they present insightful snapshots that demonstrate that future constructions offer exciting insights into processes of language change and grammaticalization, and can potentially contribute to a better understanding of the cognitive and sociocultural underpinnings of these processes.


\section*{Acknowledgments}
We are grateful to the Fritz Thyssen Foundation for funding the workshop on which this volume is based. We are also grateful to the editors of the ``Advances in Historical Linguistics" book series for the good collaboration, to the colleagues who agreed to peer-review the papers in the present volume, as well as to the team at Language Science Press and the numerous volunteer proofreaders. Last but not at all least, we would like to thank Jiabin Yao for his invaluable help in typesetting the present volume.


\sloppy
\printbibliography[heading=subbibliography,notkeyword=this]
\end{document}
